% !TeX root = ../main.tex
%%%%%%%%%%%%%%%%%%%%%%%%%%   4   %%%%%%%%%%%%%%%%%%%%%%%%%%%%%%
\section{存在问题、困难及其解决方案}

当前的主要困难是，需要在有限时间内深入理解神经网络与非线性滤波相关内容，并将其沉淀为准确、系统且可复用的知识库资料。针对学习与整理耗时较长、实际进度存在不确定性的问题，拟采取以下措施：

% 与其他进度说明保持一致：总述加编号列表，不设置 4.1 子节。
\begin{enumerate}[label=(\arabic*),leftmargin=2.5em,topsep=3pt,itemsep=6pt,parsep=0pt]
  \item \textbf{精读文献与独立思考\ }围绕核心论文反复阅读、独立推导，重点核对模型假设、数学原理与算法逻辑，对尚未理解的问题持续追踪，逐步深化认识。
  \item \textbf{结构化整理与知识复用\ }按研究问题、模型假设、推导过程和算法步骤整理知识条目，保留文献出处及适用条件，通过复核与修订形成可供后续研究使用的资料。
  \item \textbf{量化记录与进度调整\ }坚持每日制定计划和总结复盘，记录有效工作时长、疲劳程度自评及可核查成果。依据实际完成情况调整任务优先级与工作量，为难点理解和反复验证预留时间，合理估计阶段进度。
\end{enumerate}
